\chapter{Модулярно-графовое происхождение проектора копий}
\label{ch:v7-modular-copy-projector-origin-gate}

\section{Вход после копийного запрета}

Аффинно-Hodge-действие сохраняет полный $U(2)$ одинаковых копий. Проверим,
можно ли получить второй оператор не из ручной матрицы
$\operatorname{diag}(1,-1)$, а из максимального бинарного графа всех рёбер,
разрешённых строгим условием первого порядка.

Пусть $A_{\max}$ --- немеченая матрица смежности полного графа из
четырнадцати рёбер. Две физически совпадающие пары имеют одинаковые
соседства:
\begin{equation}
 N(L_L)=N(Y_L)=\{e_R,X_R,Y_R\},
 \qquad
 N(e_R)=N(X_R)=\{L_L,X_L,Y_L\}.
 \label{eq:v7-modular-copy-twin-neighbourhoods}
\end{equation}
Поэтому их разности лежат в ядре $A_{\max}$.

\section{Оператор двухшагового обхода}

Обозначим через $q_L,q_e$ проекторы на две изотипические плоскости
$\operatorname{span}\{L_L,Y_L\}$ и
$\operatorname{span}\{e_R,X_R\}$. Они определены типами представления, а не
именами старой и новой вершины. Прямой подсчёт двухшаговых путей даёт в обеих
плоскостях
\begin{equation}
 q_i A_{\max}^2 q_i
 =3\begin{pmatrix}1&1\\1&1\end{pmatrix},
 \qquad i\in\{L,e\}.
 \label{eq:v7-modular-copy-compressed-square}
\end{equation}
Следовательно, коэффициент-свободно восстанавливается обмен
\begin{equation}
 S_i=\frac13q_iA_{\max}^2q_i-q_i
 =\begin{pmatrix}0&1\\1&0\end{pmatrix},
 \qquad S_i^2=q_i,
 \label{eq:v7-modular-copy-swap}
\end{equation}
и его проекторы
\begin{equation}
 P_{i,\pm}=\frac12(q_i\pm S_i).
 \label{eq:v7-modular-copy-parity-projectors}
\end{equation}
Они не выбирают вершину из пары. Они выбирают канонические чётную и нечётную
линии. Поскольку физическая алгебра действует на совпадающих копиях
одинаково, $S_i$ коммутирует с её представлением, но не является центральным
в полной матричной алгебре родителя.

\section{Модулярный вес снимает плоскую орбиту}

На одной копийной плоскости положим $h=S$ и
\begin{equation}
 \rho_\beta=\frac{e^{-\beta h\otimes I_3}}
 {\Tr e^{-\beta h\otimes I_3}},
 \qquad \beta>0.
 \label{eq:v7-modular-copy-state}
\end{equation}
Для Hodge-вакуума на прежней плоской орбите
\begin{equation}
 M_\theta=(\cos\theta\,I_3\ \ \sin\theta\,I_3),
 \qquad M_\theta M_\theta^\dagger=I_3,
 \label{eq:v7-modular-copy-orbit}
\end{equation}
обычное действие постоянно. Но модулярно декорированный четвёртый момент
правого блока равен
\begin{equation}
 W_\beta(\theta)
 =\Tr\!\left[\rho_\beta(M_\theta^\dagger M_\theta)^2\right]
 =\frac12\left(1-\tanh\beta\,\sin2\theta\right).
 \label{eq:v7-modular-copy-weighted-moment}
\end{equation}
Он имеет единственный минимум по проекторной орбите при
\begin{equation}
 \theta=\frac\pi4\pmod\pi,
 \qquad
 W_{\min}=\frac1{1+e^{2\beta}},
 \qquad
 W_{\max}=\frac{e^{2\beta}}{1+e^{2\beta}},
 \label{eq:v7-modular-copy-extrema}
\end{equation}
а поперечный гессиан положителен:
\begin{equation}
 W_\beta''\!\left(\frac\pi4\right)=2\tanh\beta>0.
 \label{eq:v7-modular-copy-hessian}
\end{equation}
Тяжёлое пространство выбирается чётным, а лёгкое ядро --- нечётным. Общее
обращение $h\mapsto-h$ меняет их местами и является тем же двухзначным
обращением ориентации, которое уже присутствовало в модулярном родителе
Тома V.

\section{Почему это ещё не выбирает старые шесть рёбер}

Полный граф сохраняет независимые обмены двух пар. Его четырнадцать рёбер
распадаются на девять орбит группы $S_2\times S_2$. В частности,
\begin{equation}
 \{L_L-e_R,L_L-X_R,Y_L-e_R,Y_L-X_R\}
 \label{eq:v7-modular-copy-four-edge-orbit}
\end{equation}
является одной орбитой. Прежний целевой набор содержит только одно ребро из
этой четвёрки и потому не инвариантен. Значит, новый проектор не восстанавливает
старое/новое разбиение и не оправдывает пять нулевых блоков в прежнем базисе.
Он предлагает другое физическое чтение: лёгкий $H_{15}$ должен быть
нечётной комбинацией двух копий.

\section{Архитектурная граница}

$A_{\max}$ не является ранее выведенным оператором Дирака. Он получается,
если принять дополнительный принцип: все рёбра, разрешённые первым порядком,
образуют универсальный бинарный инцидентный носитель с единичными весами до
появления динамических матриц. Этот принцип не содержит непрерывного
параметра, но всё равно является новым родительским утверждением.

Без вывода $A_{\max}$ из универсальной дифференциальной оболочки,
условного ожидания или одного действия модулярный вес остаётся условным.
Подстановка фактического $D_F$ вместо $A_{\max}$ была бы круговой: его
двухшаговый оператор уже зависит от матриц, ориентацию которых требуется
вывести.

\begin{theorem}[Условный модулярный селектор копий]
Если максимальный бинарный первопорядковый граф является каноническим
додинамическим носителем, то его двухшаговый оператор коэффициент-свободно
выводит обмены $S_L,S_e$ и проекторы чётности. Соответствующее модулярное
состояние снимает непрерывную $U(2)$-орбиту Hodge-ядер и оставляет две
ориентации, связанные общим обращением. Однако происхождение самого
$A_{\max}$ и физическая редукция девяти орбит рёбер ещё не выведены.
\end{theorem}

\begin{nogocriterion}[Запрет подмены универсальной инцидентности]
Нельзя отождествлять $A_{\max}$ с физическим $D_F$ или назначать всем
разрешённым блокам единичные матрицы без отдельного родительского принципа.
Нельзя также истолковывать нечётную комбинацию как старую копию до проверки
её калибровочных токов, масс и Real-завершения.
\end{nogocriterion}

\begin{protocol}
\begin{enumerate}
 \item точные графовые близнецы: \textbf{найдены};
 \item обмен $S$ из двухшаговой инцидентности: \textbf{выведен условно};
 \item ручная метка старой вершины: \textbf{не используется};
 \item непрерывный $U(2)$ ядер: \textbf{снят до двух ориентаций};
 \item поперечный гессиан: \textbf{положителен};
 \item прежние шесть рёбер: \textbf{не восстановлены};
 \item максимальная бинарная инцидентность из одного родителя:
 \textbf{не выведена};
 \item итог: \textbf{условно положительный селектор чётности};
 \item следующий гейт: \textbf{допустимость универсального инцидентного
 родителя на полном физическом носителе}.
\end{enumerate}
\end{protocol}

Машинный сертификат сохранён в
\path{s2t_v7_modular_copy_projector_origin_gate_results.json}.

\begin{thebibliography}{9}
\bibitem{v7modcopy-modular}
F.~Ciolli, F.~Fidaleo,
\emph{Modular spectral triples and deformed Fredholm modules},
arXiv:2206.14762.
\bibitem{v7modcopy-twins}
H.~Monterde,
\emph{Strong cospectrality and twin vertices in weighted graphs},
arXiv:2111.01265.
\end{thebibliography}